% Source: https://artofproblemsolving.com/community/c3h1938343_catch_in_coordinate_plane_jhmmc_7th_grade_2019
% Stable identity: aops-topic-1938343-q01
Unfold the path: each time the ball bounces, reflect the rest of the path in the wall. Between walls the ball always moves from one wall to the other, so the bounces alternate between $y=1$ and $y=-1$. If the first bounce is on $y=1$, the unfolded walls are the lines $y=1$, $y=3$, $y=5$, and DIO's image after three reflections is $D'=(5,6)$; if it is on $y=-1$, the image is $(5,-6)$. The unfolded path is a straight segment of the same length as the real path.

\begin{asy}
usepackage("amsmath");
size(6cm);
defaultpen(fontsize(10pt));
pair J=(0,0),Dd=(5,0);
fill((0,-1)--(5.6,-1)--(5.6,1)--(0,1)--cycle,palegray);
draw((0,1)--(5.6,1),linewidth(1.1));draw((0,-1)--(5.6,-1),linewidth(1.1));
draw((0,3)--(5.6,3),dotted);draw((0,5)--(5.6,5),dotted);
pair P1=(5/6,1),P2=(2.5,-1),P3=(25/6,1);
draw(J--P1--P2--P3--Dd,blue+linewidth(1));
draw(J--(5,6),dashed);
dot(J);dot(Dd);dot((5,6));
label("$J(0,0)$",J,W);label("$D(5,0)$",Dd,SE);label("$D'(5,6)$",(5,6),E);
label("$y=1$",(5.6,1),E);label("$y=-1$",(5.6,-1),E);label("$y=3$",(5.6,3),E);label("$y=5$",(5.6,5),E);
label("\mbox{}",(min(currentpicture,user=true)+max(currentpicture,user=true))/2,black); // keeps the text after the figure black when LaTeX typesets the labels inline
\end{asy}

Hence every such path has length $\sqrt{5^2+6^2}=\sqrt{61}$, and this value is attained: the ball bounces at $\left(\dfrac{5}{6},1\right)$, $\left(\dfrac{5}{2},-1\right)$ and $\left(\dfrac{25}{6},1\right)$ (blue). The minimum distance is $\boxed{\sqrt{61}}$.
